一般 Yang--Mills 场强的非线性项固定三、四规范玻色子自相互作用。电弱规范扇区使用 $SU(2)_L$ 的结构常数 $\epsilon^{IJK}$；先在 $\mathcal W^I$ 实基中写出这些自相互作用，再用电弱混合矩阵旋转到 $W^\pm,A,Z$ 质量本征基。$SU(2)_L$ 场强为
\begin{equation}
 \mathcal W^I_{\mu\nu}
 =\partial_\mu\mathcal W^I_\nu-\partial_\nu\mathcal W^I_\mu
 -g\epsilon^{IJK}\mathcal W^J_\mu\mathcal W^K_\nu,
 \label{eq:sm-su2-fieldstrength-massrule-dp11}
\end{equation}
并定义线性旋度
\begin{equation*}
 w^I_{\mu\nu}\equiv\partial_\mu\mathcal W^I_\nu-\partial_\nu\mathcal W^I_\mu.
\end{equation*}
于是规范动能的二、三、四次项可写成
\begin{align}
 \frac{\Lag_W}{\hbar c}
 ={}&-\frac14w^I_{\mu\nu}w^{I\mu\nu}
 +\frac g2\epsilon^{IJK}w^I_{\mu\nu}\mathcal W^{J\mu}\mathcal W^{K\nu}
 \nonumber\\
 &-\frac{g^2}{4}\epsilon^{IJK}\epsilon^{ILM}
 \mathcal W^J_\mu\mathcal W^K_\nu
 \mathcal W^{L\mu}\mathcal W^{M\nu}.
 \label{eq:sm-su2-kinetic-expanded-dp11}
\end{align}
\eqnrefs{eq:qcd-three-gluon-vertex-dp10} 给出一般 Yang--Mills 三点张量；作 $f^{abc}\to\epsilon^{IJK}$、$g_s\to g$ 的群与耦合替换后，三条物理四动量 $p,q,r$ 全部流入并满足 $p+q+r=0$ 时，$SU(2)$ 三点顶角可直接写成
\begin{equation}
 \boxed{
 \Gamma^{abc}_{\mu\nu\rho}(p,q,r)
 =-g\epsilon^{abc}
 \left[
 \eta_{\mu\nu}(p-q)_\rho
 +\eta_{\nu\rho}(q-r)_\mu
 +\eta_{\rho\mu}(r-p)_\nu
 \right].}
 \label{eq:sm-yangmills-three-vertex-dp11}
\end{equation}
其中 $p,q,r$ 均采用\eqnrefs{eq:reduced-si-physical-p-r10} 的物理四动量，量纲为动量；约化图规则已经把 Fourier 变换和外腿归一化中成对出现的 $\hbar$ 因子统一提出。

质量本征基由
\begin{equation*}
 \mathcal W^\pm=\frac{\mathcal W^1\mp\ii\mathcal W^2}{\sqrt2},
 \qquad
 \mathcal W^3=s_WA+c_W\mathcal Z
\end{equation*}
给出。只有 $W^+W^-\mathcal W^3$ 组合具有非零 $SU(2)$ 结构常数，因此中性场旋转立即固定两个三点耦合：
\begin{equation}
 \boxed{
 g_{WW\gamma}=g_{\rm em},
 \qquad
 g_{WWZ}=gc_W.}
 \label{eq:sm-wwv-couplings-dp11}
\end{equation}
若 $W^+(p_+,\mu),W^-(p_-,\nu),V(k,\rho)$ 全部流入，$p_++p_-+k=0$，则
\begin{equation}
 \boxed{
 \Gamma^{WWV}_{\mu\nu\rho}
 =+\ii g_{WWV}
 \left[
 \eta_{\mu\nu}(p_+-p_-)_\rho
 +\eta_{\nu\rho}(p_--k)_\mu
 +\eta_{\rho\mu}(k-p_+)_\nu
 \right].}
 \label{eq:sm-wwv-vertex-dp11}
\end{equation}
它的洛伦兹张量结构由一般 Yang--Mills 三点顶角继承，复因子与耦合则由带电质量本征基变换固定。

电磁规范不变性要求光子纵向分量不能改变物理振幅。为把这一约束写成可直接核验的代数式，定义自由有质量矢量场的横向逆核部分
\begin{equation*}
 \Pi_{\mu\nu}(p)\equiv
 p^2\eta_{\mu\nu}-p_\mu p_\nu-m_W^2c^2\eta_{\mu\nu}.
\end{equation*}
由\eqnrefs{eq:sm-wwv-vertex-dp11} 可得
\begin{equation}
 \boxed{
 k^\rho\Gamma^{WW\gamma}_{\mu\nu\rho}
 =+\ii g_{\rm em}
 \left[\Pi_{\mu\nu}(p_-)-\Pi_{\mu\nu}(p_+)\right].}
 \label{eq:sm-wwgamma-ward-dp41}
\end{equation}
对在壳外腿 $p_\pm^2=m_W^2c^2$ 且 $p_\pm\cdot\varepsilon_\pm=0$，右端夹在物理偏振之间为零，因此把外部光子偏振替换为 $k^\rho$ 后振幅消失。

四规范玻色子顶角同样由\eqnrefs{eq:sm-su2-kinetic-expanded-dp11} 的四次项固定，并通过同一基变换产生 $W^+W^-\gamma\gamma$、$W^+W^-\gamma Z$、$W^+W^-ZZ$ 与 $W^+W^-W^+W^-$。树级不存在 $\gamma\gamma\gamma$、$ZZZ$ 或 $\gamma\gamma Z$：三个纯中性外腿都来自同一个 Cartan 方向 $\mathcal W^3$ 与阿贝尔场 $\mathcal B$，相应结构常数为零。

\begin{figure}[htbp]
\centering
\begin{adjustbox}{max width=.96\linewidth}\begin{tikzpicture}[>=Latex,node distance=1.35cm and 1.75cm,line label/.style={fill=white,fill opacity=.96,text opacity=1,inner sep=1.15pt}]
\node[draw,rounded corners,align=center] (act) {gauge-invariant action\\$G_{\rm SM}$ representations};
\node[draw,rounded corners,right=2.3cm of act,align=center] (vev) {choose Higgs vacuum\\diagonalize mass matrices};
\node[draw,rounded corners,right=of vev,align=center] (phys) {physical mass-basis vertices\\fermions, $W,Z,A,h$};
\node[draw,rounded corners,below=of vev,align=center] (gf) {$R_\xi$ gauge fixing\\uses $m_W,m_Z$};
\node[draw,rounded corners,right=of gf,align=center] (unphys) {Goldstone / ghost\\propagators and vertices};
\node[draw,rounded corners,right=of phys,align=center] (rules) {complete mass-basis\\Feynman rules};
\draw[->] (act)--node[line label,midway,above=1pt,font=\scriptsize,align=center]{EWSB +\\basis change}(vev);
\draw[->] (vev)--(phys);
\draw[->] (vev)--(gf);
\draw[->] (gf)--(unphys);
\draw[->] (phys)--(rules);
\draw[->] (unphys.east)--(rules.south west);
\end{tikzpicture}\end{adjustbox}
\caption{标准模型规范顶角的依赖关系。规范玻色子自相互作用由同一个 Yang--Mills 场强固定；$R_\xi$ 规范固定非物理内部自由度的传播结构，而物理耦合常数由规范不变作用量决定。费米子质量基中的味结构在 Yukawa 对角化后再进入相应顶角。}
\label{fig:sm-rule-dependency-dp11}
\end{figure}
